% ----------------------------------------------
% Capítol 8
% ----------------------------------------------
\chapter{El Gran Cisma i l'Herència Esquerdada}

\epigraf{Allò que amb gran esforç s'aixeca, amb la deixadesa s'ensorra.}{Sèneca}

\lettrine[loversize=0.2, lines=2]{E}{l} temps inexorable finalment va reclamar l'Eric. Després d'anys de servei implacable, l'Homogeneïtzador va conquerir els darrers crèdits i va abandonar la Universitat, posant fi a l'era dels fundadors. La custòdia del OneDrive va recaure sobre el Consell i el jove hereu, en David, just en el moment en què la institució desfermava el pitjor dels cataclismes acadèmics: la imposició d'un nou pla d'estudis. El mapa de la física es va esmicolar. Assignatures històriques van ser substituïdes per noves i críptiques matèries docents per a les quals no existia ni un sol PDF. De la nit al dia, una gran part del repositori va quedar totalment obsoleta.

Les carpetes de primer, abans fèrtils, es van convertir en un erm. Però la veritable tragèdia no va ser el buit de material, sinó l'apatia de la nova generació. Els neòfits van arribar al Drive com a simples consumidors: descarregaven allò poc que encara servia, exhaurien els recursos i marxaven sense aportar absolutament res a canvi. Aquesta actitud consumista va indignar profundament el Consell, i molt especialment, en David. Ell, que havia lluitat fins a l'extenuació per imposar la puresa del LaTeX i el rigor en els documents, observava amb amargor com la seva voluntat ferma no havia estat heretada. La mandra s'havia apoderat de les aules; el seu esforç semblava caure en el buit davant d'estudiants que es negaven a escriure codi o a perpetuar la cadena del coneixement.

Tanmateix, quan la sequera semblava definitiva, va emergir un petit reducte de resistència encapçalat per un estudiant anomenat Biel. Juntament amb un grapat d'ànimes afins, es van negar a permetre l'esfondrament total de la biblioteca. Tot i que no van arribar a generar apunts ni veritable contingut des de zero, van assumir la crucial tasca de rescatar, recopilar i penjar els exàmens de les noves assignatures. En presenciar la determinació d'aquest jove per mantenir actiu l'arxiu, en David va voler veure-hi el reflex de la seva pròpia empenta. Deixant enrere el seu desencís, va decidir adoptar oficialment en Biel sota la seva tutela directa, amb l'esperança d'instruir-lo en els sagrats dogmes de la religió del LaTeX. Però el destí és irònic i cruel. En lloc de seguir el camí de la puresa tipogràfica, el nou pupil guardava un cop amarg per al seu mestre: acabaria renegant de l'elegància del codi per acabar abraçant el pragmatisme fosc, mundà i herètic del Microsoft Word.

% ----------------------------------------------
% Capítol 9
% ----------------------------------------------
\chapter{Els Nous Apòstols}

\epigraf{Si l'equació és bella, és molt probable que sigui correcta.}{Paul Dirac}

\lettrine[loversize=0.2, lines=2]{M}{algrat} l'amarga heretgia d'en Biel i la desídia d'una gran part de les noves fornades, el destí del OneDrive no estava condemnat a la foscor. A mesura que en David i els seus fidels companys avançaven inexorablement cap als cursos superiors, no deixaven al seu pas carpetes buides ni desolació. Ben al contrari, allà per on passaven, sembraven la llum. El seu avenç per la carrera es va convertir en una autèntica croada per l'estètica i l'ordre, deixant com a llegat un repositori cada cop més ple, on tots els documents resplendien sota els inquebrantables dogmes de la religió del LaTeX.

Va ser en ple apogeu d'aquesta fe quan la providència va atorgar al OneDrive l'arribada de nous aliats. D'entre les bases va sorgir en Pol, un escriba d'una capacitat gairebé sobrehumana que va desfermar una producció titànica en pur codi. Pàgines i pàgines de demostracions complexes es van compilar amb una bellesa tipogràfica que fregava la perfecció divina. Paral·lelament, del món híbrid del doble grau en va emergir l'Héctor, un «fisquim» erudit que va assumir la missió d'omplir el buit de la seva estirp, redactant amb precisió les solucions d'aquelles matèries frontereres que els físics purs no gosaven trepitjar.

Tanmateix, una ombra d'ironia tràgica planava sobre tota aquesta magna obra. A causa de l'implacable canvi del pla docent, la geometria de segon curs s'havia deformat irreversiblement. L'esforç titànic abocat en aquells apunts i solucions perfectes no seria del tot útil per a les generacions vinents; les matèries canviaven, el temari es movia i el context desapareixia.

Per què llavors tant de sacrifici? Per què compilar milers de línies de codi si l'arxiu perdia la seva utilitat immediata?

La resposta raïa en el cor mateix de la seva fe. Per als apòstols del LaTeX, escriure ja no era una simple eina de supervivència per aprovar exàmens; s'havia convertit en un fi en si mateix. Ho feien per amor a la física, sabent que les lleis de l'univers són immutables encara que la burocràcia canviï els plans d'estudis. Escrivien per purgar el caos, per pura devoció a la bellesa, i per erigir un monument a la seva pròpia lluita. Aquells PDF no eren simples apunts: eren un manifest, una prova eterna que ells havien passat per allà, havien entès el cosmos, i s'havien negat a deixar-lo sense endreçar.




















